At $s=1$ and $0<c<1$ the loop closure is irreducible and positive recurrent
(Proposition~\ref{prop:doubling_phase}), so $\lambda$ exists and is positive at every state,
and $\varepsilon_{c,1}$ takes its values in $(0,c/2]$ with $c/2<1$. Hence
Proposition~\ref{prop:doubling_cut} gives the cut balance at every $m>d$,
Theorem~\ref{theo:doubling_decay} the two-sided bound \eqref{eq:doubling_decay},
Corollary~\ref{cor:doubling_tail} the ratio bounds \eqref{eq:doubling_tailratio} together with
$L(m)\rightarrow0$, and Lemma~\ref{lem:doubling_percut} applies in its first case at every such
cut. Choose $m_2>d$ beyond which $1-L(m)\geq\tfrac12$. There \eqref{eq:doubling_step3} bounds
the Rayleigh quotient of the centred tail indicator by $\sqrt{4\lambda_m/L(m)}$, and the second
pair of \eqref{eq:doubling_tailratio} turns that into $2m^{-1/2}\sqrt{c_2(p_*-1)/c_1}$, which is
\eqref{eq:doubling_explicit} with $c_7=\tfrac12\sqrt{c_1/(c_2(p_*-1))}$.
